\section{Problem Formulation}

\subsection{In-Forward Window Acquisition Contract}

Let a dense candidate window be
$\VideoWin=\{v_t\}_{t=1}^{\Tdense}$ with original dense-time coordinate
$\taucoord_t=t$. The main budget setting uses $\Tdense=768$ and
$\Kmain=384$. DUCA inserts an acquisition plugin in the detector forward path:
\begin{align}
    x_t &= \philo(v_t), \\
    \Ssel &= \Acquirer_{\thet}(x_{1:\Tdense})
          = \{s_1 < \cdots < s_{\Ksel}\}, \qquad
          \Ksel \le \Kmain .
    \label{eq:online_policy}
\end{align}
The detector receives only $\VideoWin_{\Ssel}$, not the full sequence. The
selected set stores original dense indices, so $\taucoord_{s_i}$ remains the
coordinate used for assignment, decoding, and evaluation. We call this setting
\emph{in-forward} or \emph{window-online}: the selection is computed from the
current detector window during the forward path and is not loaded from an
offline decision record. The setting does not require causal streaming or
prefix-only outputs.

\subsection{Observation and Compute Boundary}

The object selected by DUCA depends on where the plugin is inserted. In a
pre-backbone setting, $v_t$ is a raw clip or frame group and acquisition reduces
the cost of both the backbone and detector head. In a feature-stream setting,
$v_t$ is a precomputed temporal descriptor and acquisition reduces the cost of
the detector head and subsequent temporal processing. We therefore report
compute with separate terms for descriptor extraction, zero-shot scoring,
selector inference, backbone processing, detector head, and post-processing:
\begin{align}
    C_{\mathrm{total}}
    &=
    C_{\mathrm{desc}} + C_{\mathrm{score}} + C_{\mathrm{sel}}
    \notag\\
    &\quad
    + C_{\mathrm{backbone}}(\Ksel)
    + C_{\mathrm{head}}(\Ksel)
    + C_{\mathrm{post}} .
    \label{eq:compute_boundary}
\end{align}
This boundary prevents a frozen-feature experiment from being presented as a
full video-backbone saving.

\subsection{Teacher-Free Inference Boundary}

At validation and deployment, $x_t$ may contain only deploy-visible information:
zero-shot actionness, motion or feature energy, low-resolution descriptors, and
optional learned probes. The online decision cannot use ground truth, oracle
boundaries, teacher utility, detector raw predictions, prediction caches, or a
precomputed decision ledger. We write this as
\begin{align}
    z^{\mathrm{deploy}}_t
    &=
    (\pzs(t), e^{\mathrm{mot}}_t, e^{\mathrm{feat}}_t, p^{\mathrm{probe}}_t),
    \notag\\
    \Ssel
    &=
    \Acquirer_{\thet}(z^{\mathrm{deploy}}_{1:\Tdense}).
    \label{eq:deploy_boundary}
\end{align}
An audit ledger may record the decision after it is made, but it is not an
input to \cref{eq:online_policy}.

\subsection{Detector-Utility Target}

The selector should predict detector utility rather than actionness alone. For
training windows, a dense detector or frozen teacher defines positive detector
points $\mathcal{P}^{+}$. Each point $p$ has classification loss
$\ell^{\mathrm{cls}}_p$, localization loss $\ell^{\mathrm{loc}}_p$, ranking or
hard-negative weight $r_p$, and temporal responsibility kernel $a_{p,t}$. A
train-only utility target is
\begin{align}
    d_t &=
      \sum_{p\in\mathcal{P}^{+}}
      a_{p,t}
      (\ell^{\mathrm{cls}}_p+\ell^{\mathrm{loc}}_p+r_p),
      \notag\\
    u^{\star}_t &=
    \mathrm{norm}\!\left(
      \alpha y^{\mathrm{bnd}}_t
      + \beta y^{\mathrm{act}}_t
      + \gamma d_t
    \right).
    \label{eq:utility_target}
\end{align}
The target supervises acquisition on the training split only. It is discarded
before validation and deployment.

\subsection{Sparse Temporal Grid}

Hard selection creates a selected axis $i=1,\ldots,\Ksel$, but selected rank is
not physical time. DUCA therefore passes the detector a sparse temporal grid
\begin{equation}
    \Grid_{\Ssel}
    =
    \{(i, s_i, \taucoord_{s_i}, \ell_i, r_i)\}_{i=1}^{\Ksel},
    \label{eq:sparse_grid}
\end{equation}
where $\ell_i$ and $r_i$ are the left and right original-time support bounds of
the selected cell. With $T_{\mathrm{valid}}$ denoting the valid prefix length,
we use midpoint cells
\begin{align}
    \ell_i &=
    \begin{cases}
    1, & i=1,\\
    \frac{1}{2}(\taucoord_{s_{i-1}}+\taucoord_{s_i}), & i>1,
    \end{cases}
    \notag\\
    r_i &=
    \begin{cases}
    T_{\mathrm{valid}}, & i=\Ksel,\\
    \frac{1}{2}(\taucoord_{s_i}+\taucoord_{s_{i+1}}), & i<\Ksel.
    \end{cases}
    \label{eq:grid_cells}
\end{align}
Detector point assignment, regression range checks, offset targets, decoded
segments, and NMS are all expressed in original-time units. For a ground-truth
segment $(b^s,b^e)$ assigned to selected point $i$ with center
$c_i=\taucoord_{s_i}$, the offset target is
\begin{equation}
    \delta_i^\star=(c_i-b^s,\ b^e-c_i),
    \qquad
    \hat b^s=c_i-\hat\delta_i^L,\quad
    \hat b^e=c_i+\hat\delta_i^R .
    \label{eq:original_time_decode}
\end{equation}
If the detector predicts on the selected axis, the proposal must be remapped
through $\Grid_{\Ssel}$ before mAP is computed. This grid is the key object for
high-IoU credibility.
